\chapter{The Sampling Block}
\label{ch:sampling-block}

\newthought{The \yamlkey{Sampling} block} is the heart of a task card. It says \emph{how} points
are chosen, \emph{which} parameters vary, and \emph{how good} each point is. This chapter gives
the overall shape; the two biggest pieces---variables and the likelihood---get their own
chapters next, and each sampler is documented in Part~\ref{part:samplers}.

\section{The four parts}
\label{sec:sampling-parts}

\begin{yamlwin}
Sampling:
  Method: "Random"      # which sampler to use
  Variables: []         # the parameters to scan
  Bounds: {}            # sampler-specific settings (optional)
  LogLikelihood: []     # how to score each point
\end{yamlwin}

\begin{description}[style=nextline, leftmargin=1.2em]
  \item[\yamlkey{Method}] The name of the sampler, e.g. \sampler{Random}, \sampler{Bridson},
        \sampler{MCMC}, \sampler{Dynesty}. Choosing one is covered in
        Chapter~\ref{ch:choosing-sampler}.
  \item[\yamlkey{Variables}] The list of parameters to scan and the distribution each is drawn
        from. Full reference in Chapter~\ref{ch:variables}.
  \item[\yamlkey{Bounds}] Settings specific to the chosen \yamlkey{Method} (for example, the
        number of \textsc{mcmc} chains, or a sampler's iteration count). Each sampler's chapter
        lists its keys. Simple samplers need none.
  \item[\yamlkey{LogLikelihood}] One or more named expressions that score a point. Full
        reference in Chapter~\ref{ch:loglikelihood}.
\end{description}

\begin{jnote}
Each sampler reads its own settings from \yamlkey{Bounds} (and a few read a method-named block
such as \yamlkey{CSV} for the \sampler{CSV} replay method). Treat \yamlkey{Bounds} as ``the
dials for this sampler'' and look up the exact dials in the sampler's chapter.
\end{jnote}

\section{A worked example}
\label{sec:sampling-example}

Here is a complete \yamlkey{Sampling} block for an \sampler{MCMC} run over two variables:

\begin{yamlwin}
Sampling:
  Method: "MCMC"
  Variables:
    - name: x
      distribution: {type: Flat, parameters: {min: 0, max: 5}}
    - name: y
      distribution: {type: Flat, parameters: {min: 0, max: 5}}
  Bounds:
    num_chains: 8
    num_iters: 600
    proposal_scale: 0.06
  LogLikelihood:
    - {name: "LogL_Z", expression: "LogGauss(z, 100, 10)"}
\end{yamlwin}

\noindent Read it as a sentence: ``Run \sampler{MCMC} (with 8 chains of 600 iterations) over flat
variables \yamlkey{x} and \yamlkey{y} on $[0,5]$, scoring each point by how close the observable
\yamlkey{z} is to $100 \pm 10$.''

\section{Optional: pre-selecting points}
\label{sec:sampling-selection}

Some samplers accept an optional \yamlkey{selection}: a true/false condition that a point must
satisfy to be accepted. It is written as a symbolic expression
(Chapter~\ref{ch:symbols}):

\begin{yamlwin}
Sampling:
  selection: "(x > 0) & (y < 4)"
\end{yamlwin}

\noindent Use it to skip obviously uninteresting regions before the (often expensive) workflow
runs. The expression must evaluate to a strict boolean.

\section{Where to go next}
\label{sec:sampling-next}

\begin{itemize}
  \item To define the variables: Chapter~\ref{ch:variables}.
  \item To define the score: Chapter~\ref{ch:loglikelihood}.
  \item To choose and tune a sampler: Chapter~\ref{ch:choosing-sampler} and the per-method
        chapters in Part~\ref{part:samplers}.
\end{itemize}
